\chapter{Zaključak i budući rad}

Ovaj rad je pokazao da se prihvatljivo tačan, transparentan i \textbf{objašnjiv} sistem za rutiranje teretnih vozila može izgraditi nad potpuno otvorenim geografskim podacima (OpenStreetMap) i otvorenim routing engine-om (Valhalla), uz sopstveni algoritamski doprinos koji tim alatima nedostaje. Implementiran je kompletan sistem koji obuhvata: pripremu geografskih podataka specifičnih za teretni saobraćaj (poglavlje 5), backend servis koji generiše, rangira sopstvenom funkcijom rizika i objašnjava predloženu rutu (poglavlje 6 i 7), sopstvenu, nezavisno testiranu Dijkstra/A* implementaciju nad ograničenim, ali stvarnim podacima (poglavlje 6.4, evaluirana u poglavlju 9.2), distribuiranu komunikaciju preko REST API-ja, RabbitMQ posrednika poruka i WebSocket veza (poglavlje 7), i mobilnu aplikaciju za dve različite korisničke uloge — vozača i dispečera (poglavlje 8). Konkretan slučaj (visinsko ograničenje kod Novih Banovaca, poglavlje 6.3) je korišćen kao stvaran, provaljen dokaz da mehanizam objašnjenja rute radi na pravom, a ne samo hipotetičkom problemu, a slučaj Radalj (poglavlje 9.3) kao konkretan primer u kome je sopstvena funkcija rizika ispravila objektivno pogrešan izbor rute.

\section{Doprinosi rada}

Glavni doprinosi rada su:

\begin{enumerate}
\item Sopstveni sloj rangiranja alternativnih ruta prema riziku specifičnom za teretni saobraćaj i preferencama vozača, nad Valhalla alternativama koje engine sam ne rangira (poglavlje 6.2).
\item Automatizovan mehanizam koji vozaču objašnjava zašto predložena ruta odstupa od geometrijski najkraćeg puta, sa geometrijskom (ne pozicionom) detekcijom mesta odstupanja, potvrđen na stvarnom slučaju (poglavlje 6.3).
\item Samostalna, testirana Dijkstra/A* implementacija nad ograničenim podgrafom OSM podataka, sa sopstvenom funkcijom cene i isključivanjem ivica/čvorova na osnovu fizičkih ograničenja vozila, evaluirana nad stvarnim, a ne samo sintetičkim podacima (poglavlje 6.4, 9.2).
\item Kompletan, iako po obimu ograničen, distribuiran sistem — REST API, asinhrona obrada porukama, WebSocket komunikacija u realnom vremenu — sa modelom uloga dispečer/vozač i mehanizmom ponude/prihvatanja ture (poglavlje 7).
\end{enumerate}

\section{Ograničenja rada}

U skladu sa realno postavljenim obimom (obrazloženim već u poglavlju 1.2), rad svesno ne obuhvata:

\begin{itemize}
\item \textbf{Punu AETR regulativu} o radnom vremenu i pauzama vozača (poglavlje 7.4) — trenutni prag od 270 minuta je pojednostavljena zamena, ne zakonski obavezujuća logika.
\item \textbf{Elevacioni (nagibni) rizik rute} — sistem nema pristup podacima o nagibu puta, pa je procena potrošnje goriva (poglavlje 6.2) relativna, ne apsolutna veličina.
\item \textbf{Offline režim rada mobilne aplikacije} — aplikacija zahteva internet konekciju za svaku operaciju; rad bez signala/na lokalnoj mreži nije podržan.
\item \textbf{Podršku za više država} — geografski podaci i pripremljen graf pokrivaju isključivo teritoriju Republike Srbije.
\item \textbf{Produkcioni nivo bezbednosti} — WebSocket \texttt{CheckOrigin} je propustljiv (poglavlje 4.3), a RabbitMQ topologija je minimalna (jedan topic exchange, bez retry/DLQ mehanizma), prihvatljivo za sistem sa jednim poznatim klijentom, ne za javni, višezakupčani servis.
\item \textbf{Kalibraciju koeficijenata funkcije rizika i funkcije cene} prema stvarnim empirijskim podacima (poglavlje 9.4) — koeficijenti su prva heuristička procena.
\item \textbf{Prostorne upite nad geometrijom rute direktno u bazi} — geometrija se čuva kao tekstualni polyline, ne PostGIS \texttt{geometry} kolona (poglavlje 7.2), iako je PostGIS ekstenzija već deo infrastrukture.
\end{itemize}

\section{Pravci budućeg rada}

Na osnovu ograničenja iz 10.2, prirodni pravci budućeg razvoja su: (1) implementacija pune AETR logike, uključujući razliku između dnevnog i nedeljnog vremena vožnje i zakonski propisanih izuzetaka; (2) uvođenje elevacionih podataka (npr. iz SRTM ili sličnog izvora) u funkciju rizika, radi realnije procene potrošnje goriva na deonicama sa izraženim nagibom; (3) offline režim mobilne aplikacije sa lokalnim kešom poslednje poznate rute i odloženim slanjem GPS pozicija po povratku konekcije; (4) proširenje geografskog obuhvata na region ili celu Evropu, što bi zahtevalo i preispitivanje da li ograničeni, u memoriji učitan skup odmarališta (poglavlje 5.3) i dalje skalira, ili treba preći na PostGIS prostorni upit nad već postojećom ekstenzijom; (5) prikupljanje stvarnih podataka (npr. iz simulacije ili pilot upotrebe) radi kalibracije koeficijenata funkcije rizika (6.2) i funkcije cene sopstvenog algoritma (6.4) regresijom ili sličnim metodom, umesto ručno pretpostavljenih vrednosti; (6) bezbednosni hardening WebSocket sloja (stroža \texttt{CheckOrigin} provera, autentifikacija na nivou pojedinačne poruke) i proizvodna RabbitMQ topologija (retry, dead-letter queue) ako bi sistem prešao iz demonstracionog u produkcioni kontekst sa nepoznatim, nepoverljivim klijentima.

\clearpage
